\chapter{Multiple Integrals}
\label{ch:multint}

A single integral adds up values along a line. A multiple integral adds up values over a plane region or a solid. In this chapter we define double and triple integrals, prove that they can be computed one variable at a time (Fubini's theorem), and prove the change-of-variables formula, which explains the factors $r$ and $\rho^2\sin\varphi$ in polar and spherical coordinates. We finish with two techniques from analysis that every interview candidate should know: the Gaussian integral and differentiation under the integral sign.

%====================================================================
\section{The double integral}
%====================================================================

Let $R=[a,b]\times[c,d]$ be a rectangle and $f:R\to\R$ a bounded function. A \emph{partition} $P$ of $R$ splits $[a,b]$ into $a=x_0<\dots<x_m=b$ and $[c,d]$ into $c=y_0<\dots<y_n=d$. This gives $mn$ small rectangles $R_{ij}=[x_{i-1},x_i]\times[y_{j-1},y_j]$ with area $\Delta A_{ij}=\Delta x_i\Delta y_j$. Let
\[
m_{ij}=\inf_{R_{ij}}f,\qquad M_{ij}=\sup_{R_{ij}}f,\qquad L(P)=\sum m_{ij}\Delta A_{ij},\qquad U(P)=\sum M_{ij}\Delta A_{ij}.
\]
$L(P)$ and $U(P)$ are the \emph{lower} and \emph{upper sums}.

\begin{definition}[Double integral]
The function $f$ is \emph{integrable} on $R$ if $\sup_PL(P)=\inf_PU(P)$. The common value is the \emph{double integral} $\displaystyle\iint_Rf\,\dd A$.
\end{definition}

If $f\ge0$, the double integral is the volume of the solid under the graph of $f$ and above $R$: the lower and upper sums are the volumes of columns that fit inside and outside the solid (Figure~\ref{fig:riemann}).

\begin{figure}[ht]
\centering
\tdplotsetmaincoords{62}{128}
\begin{tikzpicture}[tdplot_main_coords,scale=1.35]
  \def\n{4}
  \foreach \i in {0,1,2,3}{
    \foreach \j in {0,1,2,3}{
      \pgfmathsetmacro{\xc}{\i+0.5}
      \pgfmathsetmacro{\yc}{\j+0.5}
      \pgfmathsetmacro{\hh}{0.35+0.12*(\xc)+0.1*(\yc)}
      \pgfmathsetmacro{\xnext}{\i+1}
      \pgfmathsetmacro{\ynext}{\j+1}
      \fill[accent!18] (\i,\j,\hh)--(\xnext,\j,\hh)--(\xnext,\ynext,\hh)--(\i,\ynext,\hh)--cycle;
      \fill[accent!35] (\xnext,\j,0)--(\xnext,\ynext,0)--(\xnext,\ynext,\hh)--(\xnext,\j,\hh)--cycle;
      \fill[accent!28] (\i,\ynext,0)--(\xnext,\ynext,0)--(\xnext,\ynext,\hh)--(\i,\ynext,\hh)--cycle;
      \draw[accent!70!ink,thin] (\i,\j,\hh)--(\xnext,\j,\hh)--(\xnext,\ynext,\hh)--(\i,\ynext,\hh)--cycle;
      \draw[accent!70!ink,thin] (\xnext,\j,0)--(\xnext,\j,\hh) (\xnext,\ynext,0)--(\xnext,\ynext,\hh) (\i,\ynext,0)--(\i,\ynext,\hh);
    }
  }
  \draw[axisline,mid] (0,0,0)--(4.8,0,0) node[below right]{$x$};
  \draw[axisline,mid] (0,0,0)--(0,4.8,0) node[right]{$y$};
  \draw[axisline,mid] (0,0,0)--(0,0,2.4) node[above]{$z$};
\end{tikzpicture}
\caption{A Riemann sum for $\iint_Rf\,\dd A$: the solid under the surface $z=f(x,y)$ is approximated by columns over small rectangles. The sum of the column volumes converges to the double integral as the rectangles shrink.}
\label{fig:riemann}
\end{figure}

\begin{theorem}[Continuous functions are integrable]\label{thm:contint}
If $f$ is continuous on $R$, then $f$ is integrable, and $U(P)-L(P)\to0$ as the mesh of $P$ (the largest side of its rectangles) tends to $0$.
\end{theorem}

\begin{proof}
A continuous function on the compact set $R$ is uniformly continuous: for each $\eps>0$ there is $\delta>0$ such that $\abs{f(\vec p)-f(\vec q)}<\eps$ whenever $\norm{\vec p-\vec q}<\delta$. If the mesh of $P$ is smaller than $\delta/2$, each small rectangle has diameter less than $\delta$, so $M_{ij}-m_{ij}\le\eps$. Then $U(P)-L(P)=\sum(M_{ij}-m_{ij})\Delta A_{ij}\le\eps\sum\Delta A_{ij}=\eps\,\Area(R)$. Since $L(P)\le\sup L\le\inf U\le U(P)$, we get $0\le\inf U-\sup L\le\eps\Area(R)$ for every $\eps$; hence $\sup L=\inf U$.
\end{proof}

The double integral is \emph{linear} ($\iint(\alpha f+\beta g)=\alpha\iint f+\beta\iint g$), \emph{monotone} ($f\le g\Rightarrow\iint f\le\iint g$) and \emph{additive} over regions that overlap only along boundaries. For a general bounded region $D$ whose boundary has zero area, define $\iint_Df\,\dd A=\iint_Rf\chi_D\,\dd A$, where $R\supseteq D$ is a rectangle and $\chi_D$ equals $1$ on $D$ and $0$ outside. In particular $\Area(D)=\iint_D1\,\dd A$.

%====================================================================
\section{Fubini's theorem: iterated integrals}
%====================================================================

\begin{theorem}[Fubini for continuous functions]\label{thm:fubini}
If $f$ is continuous on $R=[a,b]\times[c,d]$, then
\[
\iint_Rf\,\dd A=\int_a^b\Bigl(\int_c^df(x,y)\,\dd y\Bigr)\dd x=\int_c^d\Bigl(\int_a^bf(x,y)\,\dd x\Bigr)\dd y .
\]
\end{theorem}

\begin{proof}
Let $F(x)=\int_c^df(x,y)\,\dd y$. This function is continuous in $x$ (by uniform continuity of $f$: $\abs{F(x)-F(x')}\le(d-c)\sup_y\abs{f(x,y)-f(x',y)}$). Take a partition $P$ and fix $i$. For $x\in[x_{i-1},x_i]$ and each $j$, we have $m_{ij}\le f(x,y)\le M_{ij}$ for $y\in[y_{j-1},y_j]$. Integrating over this $y$-interval and summing over $j$,
\[
\sum_jm_{ij}\Delta y_j\ \le\ F(x)\ \le\ \sum_jM_{ij}\Delta y_j .
\]
Now integrate over $x\in[x_{i-1},x_i]$ and sum over $i$:
\[
L(P)\ \le\ \int_a^bF(x)\,\dd x\ \le\ U(P).
\]
By definition, $L(P)\le\iint_Rf\le U(P)$ too. By Theorem~\ref{thm:contint}, $U(P)-L(P)\to0$ as the mesh tends to $0$. So the two numbers $\int_aF$ and $\iint_Rf$, which both lie between $L(P)$ and $U(P)$, differ by less than any $\eps$; they are equal. Exchanging the roles of $x$ and $y$ gives the second equality.
\end{proof}

\begin{intuition}
This is Cavalieri's principle: to find a volume, cut the solid into thin slices, find the area of each slice (an ordinary integral), and add the slices (another ordinary integral).
\end{intuition}

For non-rectangular regions, Fubini's theorem still holds. A region is of \emph{type I} if $D=\{(x,y):a\le x\le b,\ g_1(x)\le y\le g_2(x)\}$ with $g_1\le g_2$ continuous, and of \emph{type II} if $D=\{(x,y):c\le y\le d,\ h_1(y)\le x\le h_2(y)\}$. Then, for $f$ continuous on $D$,
\[
\iint_Df\,\dd A=\int_a^b\!\int_{g_1(x)}^{g_2(x)}f(x,y)\,\dd y\,\dd x\quad\text{(type I)},\qquad
\iint_Df\,\dd A=\int_c^d\!\int_{h_1(y)}^{h_2(y)}f(x,y)\,\dd x\,\dd y\quad\text{(type II)}.
\]
(To prove this, extend $f$ by $0$ outside $D$ to a rectangle; the extended function is discontinuous only along the graphs of $g_1,g_2$, which have zero area, and the proof above goes through with the same estimates because $U(P)-L(P)\to0$ still holds. We omit this routine modification.)

\begin{figure}[ht]
\centering
\begin{tikzpicture}[x=3.1cm,y=3.1cm]
  % left: type I strip
  \begin{scope}
    \draw[axisline,mid] (-0.05,0)--(1.15,0) node[right]{$x$};
    \draw[axisline,mid] (0,-0.05)--(0,1.15) node[above]{$y$};
    \path[regionfill] plot[domain=0:1,samples=40](\x,{sqrt(\x)}) -- plot[domain=1:0,samples=40](\x,{\x*\x}) -- cycle;
    \fill[accent2!25] (0.44,0.1936) rectangle (0.48,0.6928);
    \draw[accent2,very thick] (0.46,0.2116)--(0.46,0.6782);
    \node[accent2,font=\small] at (0.46,0.80) {$y$ from $x^2$ to $\sqrt x$};
    \node[accent,font=\small] at (0.3,0.35) {$D$};
    \node[font=\small] at (0.55,-0.1) {type I: vertical strips};
  \end{scope}
  % right: type II strip
  \begin{scope}[xshift=5.6cm]
    \draw[axisline,mid] (-0.05,0)--(1.15,0) node[right]{$x$};
    \draw[axisline,mid] (0,-0.05)--(0,1.15) node[above]{$y$};
    \path[regionfill] plot[domain=0:1,samples=40](\x,{sqrt(\x)}) -- plot[domain=1:0,samples=40](\x,{\x*\x}) -- cycle;
    \fill[accent3!25] (0.2025,0.44) rectangle (0.6708,0.48);
    \draw[accent3,very thick] (0.2025,0.46)--(0.6708,0.46);
    \node[accent3,font=\small] at (0.95,0.46) {$x$ from $y^2$ to $\sqrt y$};
    \node[font=\small] at (0.55,-0.1) {type II: horizontal strips};
  \end{scope}
\end{tikzpicture}
\caption{The region between $y=x^2$ and $y=\sqrt x$ can be described in two ways. Left: for fixed $x$, $y$ runs from $x^2$ to $\sqrt x$. Right: for fixed $y$, $x$ runs from $y^2$ to $\sqrt y$. Both give the same integral, but one order is often much easier.}
\label{fig:regions}
\end{figure}

\begin{example}[Swapping the order of integration]
Compute $\displaystyle\int_0^1\!\!\int_x^1e^{y^2}\,\dd y\,\dd x$. The inner integral has no elementary form. The region is $0\le x\le y\le1$ (a triangle), which is also $\{0\le y\le1,\ 0\le x\le y\}$. So
\[
\int_0^1\!\!\int_0^ye^{y^2}\,\dd x\,\dd y=\int_0^1ye^{y^2}\,\dd y=\tfrac12\bigl(e-1\bigr).
\]
\end{example}

\begin{example}[Area between curves]
The area of the region of Figure~\ref{fig:regions} is $\int_0^1\!\int_{x^2}^{\sqrt x}1\,\dd y\,\dd x=\int_0^1(\sqrt x-x^2)\,\dd x=\tfrac23-\tfrac13=\tfrac13$.
\end{example}

\begin{pitfall}
Fubini's theorem can fail when $f$ is not continuous (or not absolutely integrable). Example: $f(x,y)=\dfrac{x^2-y^2}{(x^2+y^2)^2}$ on $(0,1]^2$ gives $\int_0^1\!\int_0^1f\,\dd y\,\dd x=\frac\pi4$ but $\int_0^1\!\int_0^1f\,\dd x\,\dd y=-\frac\pi4$. The function is unbounded near the origin, so it is not covered by Theorem~\ref{thm:fubini}.
\end{pitfall}

%====================================================================
\section{Triple integrals}
%====================================================================

The definition and theorems carry over to $\R^3$: for a box $B=[a,b]\times[c,d]\times[e,f]$ and continuous $F$, the triple integral $\iiint_BF\,\dd V$ is the limit of sums $\sum F(\vec p_k)\Delta V_k$, and it equals any iterated integral in any order (the proof of Theorem~\ref{thm:fubini} applies to each pair of variables in turn). For a solid $E$ between two surfaces, $E=\{(x,y)\in D,\ g_1(x,y)\le z\le g_2(x,y)\}$,
\[
\iiint_EF\,\dd V=\iint_D\Bigl(\int_{g_1(x,y)}^{g_2(x,y)}F(x,y,z)\,\dd z\Bigr)\dd A .
\]
The volume of $E$ is $\iiint_E1\,\dd V$.

\begin{example}[Volume of a tetrahedron]
Let $E$ be the tetrahedron $x,y,z\ge0$, $x+y+z\le1$. Then
\[
\Vol(E)=\int_0^1\!\!\int_0^{1-x}\!\!\int_0^{1-x-y}\dd z\,\dd y\,\dd x=\int_0^1\!\!\int_0^{1-x}(1-x-y)\,\dd y\,\dd x=\int_0^1\tfrac12(1-x)^2\dd x=\tfrac16 .
\]
\end{example}

\begin{example}[Mass and centre of mass]
If a solid $E$ has density $\rho(\vec x)$, its mass is $m=\iiint_E\rho\,\dd V$ and its centre of mass is $\bar{\vec x}=\frac1m\iiint_E\vec x\,\rho\,\dd V$. Its moment of inertia about the $z$-axis is $I_z=\iiint_E(x^2+y^2)\rho\,\dd V$.
\end{example}

%====================================================================
\section{Change of variables}
%====================================================================

A difficult region or integrand often becomes simple after a change of variables. The substitution rule $\int f(x)\,\dd x=\int f(g(u))g'(u)\,\dd u$ has this generalization.

\begin{theorem}[Change of variables in a double integral]\label{thm:cov}
Let $D\subset\R^2$ be a closed bounded region whose boundary has zero area. Let $\vec\varphi:W\to\R^2$ be a $C^1$ map, defined on an open set $W\supseteq D$, that is injective on $D$ and has $\det\Jac\vec\varphi\ne0$ on $W$. Put $D'=\vec\varphi(D)$. For every continuous $f:D'\to\R$,
\[
\iint_{D'}f(x,y)\,\dd x\,\dd y=\iint_Df\bigl(\vec\varphi(u,v)\bigr)\,\bigl\lvert\det\Jac\vec\varphi(u,v)\bigr\rvert\,\dd u\,\dd v .
\]
The same statement holds in $\R^n$, with $\abs{\det\Jac\vec\varphi}$ the absolute value of the $n\times n$ Jacobian determinant.
\end{theorem}

The idea (Figure~\ref{fig:cov}): a tiny square of side $s$ in the $(u,v)$-plane is mapped almost to a parallelogram, with area $\abs{\det\Jac\vec\varphi}\,s^2$. Adding the contributions gives the formula. Here is the proof in the case $n=2$.

\begin{proof}
Write $J(\vec u)=\abs{\det\Jac\vec\varphi(\vec u)}$.

\emph{Step 1 (area of a small image).} We claim: for every $\eta>0$ there is $\delta>0$ such that every closed square $Q\subset W$ with centre $\vec u$ and side $s<\delta$ satisfies
\[
\Area\bigl(\vec\varphi(Q)\bigr)\le(1+\eta)\,J(\vec u)\,s^2 .
\]
Proof of the claim. The sets $D$ and the closure of a small neighbourhood of $D$ are compact, so $\Jac\vec\varphi$ is uniformly continuous there, and $\norm{(\Jac\vec\varphi(\vec u))^{-1}}\le K$ for some constant $K$. Let $A=\Jac\vec\varphi(\vec u)$ and $\vec g(\vec h)=\vec\varphi(\vec u+\vec h)-\vec\varphi(\vec u)-A\vec h$. Then $\mathrm{D}\vec g(\vec h)=\Jac\vec\varphi(\vec u+\vec h)-A$, whose norm is smaller than any given $\eps'$ if $\norm{\vec h}$ is small enough (uniform continuity). By the mean value inequality (Lemma~\ref{lem:mvi}), $\norm{\vec g(\vec h)}\le\eps'\norm{\vec h}$. Hence
\[
\vec\varphi(\vec u+\vec h)=\vec\varphi(\vec u)+A\bigl(\vec h+\vec e(\vec h)\bigr),\qquad\vec e(\vec h)=A^{-1}\vec g(\vec h),\quad\norm{\vec e(\vec h)}\le K\eps'\norm{\vec h}.
\]
For $\vec h$ in the square $Q-\vec u=[-s/2,s/2]^2$ we have $\norm{\vec h}\le s/\sqrt2$, so $\norm{\vec e(\vec h)}\le K\eps's/\sqrt2$: the point $\vec h+\vec e(\vec h)$ is at distance at most $K\eps's/\sqrt2$ from the square $[-s/2,s/2]^2$. Hence it lies in the square of side $s(1+\eta')$ centred at $\vec 0$, where $\eta'=\sqrt2\,K\eps'$ can be made as small as we like by choosing $\eps'$ small. Therefore $\vec\varphi(Q)$ is contained in $\vec\varphi(\vec u)+A(Q')$, where $Q'$ is a square of side $s(1+\eta')$. By Proposition~\ref{prop:detarea}, $\Area(A(Q'))=\abs{\det A}\,s^2(1+\eta')^2$. Choose $\eps'$ so small that $(1+\eta')^2\le1+\eta$. This proves the claim.

\emph{Step 2 (one inequality).} Let $f\ge0$ be continuous. Cover $D$ by a grid of closed squares $Q_k$ of side $s<\delta$ with centres $\vec u_k$ (those that meet $D$). The images $\vec\varphi(Q_k)$ cover $D'$, so
\[
\iint_{D'}f\le\sum_k\iint_{\vec\varphi(Q_k)}f\le\sum_k\Bigl(\max_{\vec\varphi(Q_k)}f\Bigr)\Area\bigl(\vec\varphi(Q_k)\bigr)\le(1+\eta)\sum_k\Bigl(\max_{\vec\varphi(Q_k)}f\Bigr)J(\vec u_k)\,s^2 .
\]
As $s\to0$, the right-hand sum tends to $(1+\eta)\iint_Df(\vec\varphi)\,J$ (it is a Riemann sum for the continuous function $(f\circ\vec\varphi)J$; squares that stick out of $D$ contribute an amount that tends to $0$ because the boundary of $D$ has zero area). Since $\eta$ is arbitrary,
\[
\iint_{D'}f\,\dd A\le\iint_D(f\circ\vec\varphi)\,J\,\dd A\qquad(f\ge0).\tag{$*$}
\]

\emph{Step 3 (the reverse inequality).} The inverse map $\vec\psi=\vec\varphi^{-1}$ is $C^1$ by the inverse function theorem (Theorem~\ref{thm:inverse}), and by the chain rule $J_{\vec\varphi}(\vec\psi(\vec x))\cdot J_{\vec\psi}(\vec x)=1$. Apply $(*)$ to $\vec\psi$ and to the function $g=(f\circ\vec\varphi)J_{\vec\varphi}\ge0$ on $D$:
\[
\iint_D(f\circ\vec\varphi)J_{\vec\varphi}\,\dd A\le\iint_{D'}g(\vec\psi(\vec x))\,J_{\vec\psi}(\vec x)\,\dd A=\iint_{D'}f(\vec x)\,J_{\vec\varphi}(\vec\psi(\vec x))J_{\vec\psi}(\vec x)\,\dd A=\iint_{D'}f\,\dd A.
\]
Together with $(*)$ this proves equality for $f\ge0$. For a general continuous $f$ write $f=f^+-f^-$ with $f^\pm\ge0$ continuous.
\end{proof}

\begin{figure}[ht]
\centering
\begin{tikzpicture}[scale=1.3]
  % left: uv-plane with small square
  \draw[axisline,mid] (-0.2,0)--(3.3,0) node[right]{$u$};
  \draw[axisline,mid] (0,-0.2)--(0,2.8) node[above]{$v$};
  \foreach \i in {0,0.5,...,3}{ \draw[mid!40,thin] (\i,0)--(\i,2.5); }
  \foreach \j in {0,0.5,...,2.5}{ \draw[mid!40,thin] (0,\j)--(3,\j); }
  \filldraw[fill=accent2!40,draw=accent2,thick] (1.5,1.0) rectangle (2.0,1.5);
  \node[accent2,font=\small] at (2.55,1.9) {area $s^2$};
  \draw[-{Stealth[length=3mm]},very thick,mid] (3.6,1.4)--(4.9,1.4) node[midway,above]{$\vec\varphi$};
  % right: curved image (polar-like map u->r in [1,3], v-> angle)
  \begin{scope}[xshift=5.2cm,yshift=-0.2cm]
    \draw[axisline,mid] (-0.2,0)--(3.6,0) node[right]{$x$};
    \draw[axisline,mid] (0,-0.2)--(0,3.4) node[above]{$y$};
    \foreach \r in {0.5,1,...,3}{ \draw[mid!40,thin] (0,\r) arc (90:0:\r); }
    \foreach \a in {0,15,...,90}{ \draw[mid!40,thin] (\a:0.5)--(\a:3); }
    \filldraw[fill=accent2!40,draw=accent2,thick] (30:1.5) arc (30:45:1.5) -- (45:2.0) arc (45:30:2.0) -- cycle;
    \node[accent2,font=\small] at (1.2,2.9) {area $\approx J\,s^2$};
  \end{scope}
\end{tikzpicture}
\caption{The map $\vec\varphi$ sends a small square (left) to a slightly curved quadrilateral (right). Because $\vec\varphi$ is almost linear on small scales, the image has area close to $\abs{\det\Jac\vec\varphi}$ times the area of the square. Here $\vec\varphi$ is the polar-coordinate map.}
\label{fig:cov}
\end{figure}

\subsection{The standard coordinate systems}

\begin{theorem}[Area and volume elements]\label{thm:elements}
\begin{enumerate}[label=(\alph*)]
  \item Polar: $x=r\cos\theta$, $y=r\sin\theta$ gives $\dd A=r\,\dd r\,\dd\theta$.
  \item Cylindrical: $x=r\cos\theta$, $y=r\sin\theta$, $z=z$ gives $\dd V=r\,\dd r\,\dd\theta\,\dd z$.
  \item Spherical: $x=\rho\sin\varphi\cos\theta$, $y=\rho\sin\varphi\sin\theta$, $z=\rho\cos\varphi$ ($\rho\ge0$, $0\le\varphi\le\pi$) gives $\dd V=\rho^2\sin\varphi\,\dd\rho\,\dd\varphi\,\dd\theta$.
\end{enumerate}
\end{theorem}

\begin{proof}
(a) was computed in Chapter~\ref{ch:diff}: $\det\Jac=r$. (b) The Jacobian is block-triangular with the polar block and a $1$ for $z$, so $\det=r$.

(c) The rows of the Jacobian matrix (derivatives with respect to $\rho,\varphi,\theta$) are
\begin{align*}
\vec r_\rho&=(\sin\varphi\cos\theta,\ \sin\varphi\sin\theta,\ \cos\varphi),\\
\vec r_\varphi&=\rho\,(\cos\varphi\cos\theta,\ \cos\varphi\sin\theta,\ -\sin\varphi),\\
\vec r_\theta&=\rho\sin\varphi\,(-\sin\theta,\ \cos\theta,\ 0).
\end{align*}
A direct computation shows that they are pairwise perpendicular: for example $\vec r_\rho\cdot\vec r_\varphi=\rho(\sin\varphi\cos\varphi\cos^2\theta+\sin\varphi\cos\varphi\sin^2\theta-\cos\varphi\sin\varphi)=0$, $\vec r_\rho\cdot\vec r_\theta=\rho\sin^2\varphi(-\cos\theta\sin\theta+\sin\theta\cos\theta)=0$, and $\vec r_\varphi\cdot\vec r_\theta=\rho^2\sin\varphi\cos\varphi(-\cos\theta\sin\theta+\sin\theta\cos\theta)=0$. Their lengths are $\norm{\vec r_\rho}=1$, $\norm{\vec r_\varphi}=\rho$, $\norm{\vec r_\theta}=\rho\sin\varphi$. The volume of the box spanned by three perpendicular vectors is the product of the lengths (Theorem~\ref{thm:triple}), so $\abs{\det\Jac}=\rho^2\sin\varphi$ (with $\sin\varphi\ge0$ on $[0,\pi]$).
\end{proof}

\begin{figure}[ht]
\centering
\begin{minipage}{0.48\textwidth}
\centering
\begin{tikzpicture}[scale=1.5]
  \draw[axisline,mid] (-0.2,0)--(3.0,0) node[right]{$x$};
  \draw[axisline,mid] (0,-0.2)--(0,3.0) node[above]{$y$};
  \draw[regionfill] (25:1.6) arc (25:55:1.6) -- (55:2.4) arc (55:25:2.4) -- cycle;
  \draw[guide] (0,0)--(25:2.6);
  \draw[guide] (0,0)--(55:2.6);
  \draw[accent2,very thick] (25:1.6)--(25:2.4);
  \node[accent2,font=\small] at (25:2.0) [below right] {$\dd r$};
  \draw[accent3,very thick] (25:2.4) arc (25:55:2.4);
  \node[accent3,font=\small] at (40:2.6) [above right] {$r\,\dd\theta$};
  \draw[accent] (0.6,0) arc (0:25:0.6);
  \node[font=\scriptsize] at (0.9,0.12) {$\theta$};
\end{tikzpicture}
\end{minipage}
\begin{minipage}{0.48\textwidth}
\centering
\tdplotsetmaincoords{66}{118}
\begin{tikzpicture}[tdplot_main_coords,scale=1.7]
  % local frame at P (rho=2, phi=60, theta=40)
  \def\Px{1.327}\def\Py{1.113}\def\Pz{1.0}
  \def\erx{0.6634}\def\ery{0.5567}\def\erz{0.5}
  \def\efx{0.383}\def\efy{0.3214}\def\efz{-0.866}
  \def\etx{-0.643}\def\ety{0.766}\def\etz{0}
  \def\L{0.8}
  \newcommand{\C}[3]{({\Px+#1*\L*\erx+#2*\L*\efx+#3*\L*\etx},{\Py+#1*\L*\ery+#2*\L*\efy+#3*\L*\ety},{\Pz+#1*\L*\erz+#2*\L*\efz+#3*\L*\etz})}
  \draw[axisline,mid] (0,0,0)--(2.4,0,0) node[right]{$x$};
  \draw[axisline,mid] (0,0,0)--(0,2.4,0) node[right]{$y$};
  \draw[axisline,mid] (0,0,0)--(0,0,2.4) node[above]{$z$};
  \draw[guide] (0,0,0)--(\Px,\Py,\Pz);
  \foreach \ib in {0,1}{\foreach \ic in {0,1}{
    \draw[accent2,very thick] \C{0}{\ib}{\ic}--\C{1}{\ib}{\ic};
    \draw[accent3,very thick] \C{\ib}{0}{\ic}--\C{\ib}{1}{\ic};
    \draw[accent4,very thick] \C{\ib}{\ic}{0}--\C{\ib}{\ic}{1};
  }}
  \node[accent2,font=\scriptsize] at \C{0.5}{0}{1.35} {$\dd\rho$};
  \node[accent3,font=\scriptsize] at \C{1.3}{0.5}{0} {$\rho\,\dd\varphi$};
  \node[accent4,font=\scriptsize] at \C{0}{1.3}{0.5} {$\rho\sin\varphi\,\dd\theta$};
\end{tikzpicture}
\end{minipage}
\caption{Volume elements. Left: in polar coordinates a small ``rectangle'' has sides $\dd r$ and $r\,\dd\theta$, so its area is $r\,\dd r\,\dd\theta$. Right: in spherical coordinates a small ``box'' has perpendicular edges $\dd\rho$, $\rho\,\dd\varphi$ and $\rho\sin\varphi\,\dd\theta$ (the radius of the circle of latitude is $\rho\sin\varphi$), so its volume is $\rho^2\sin\varphi\,\dd\rho\,\dd\varphi\,\dd\theta$.}
\label{fig:elements}
\end{figure}

\begin{example}[Volume of a ball]
For the ball of radius $R$,
\[
\Vol=\int_0^{2\pi}\!\!\int_0^\pi\!\!\int_0^R\rho^2\sin\varphi\,\dd\rho\,\dd\varphi\,\dd\theta=\frac{R^3}3\cdot\bigl[-\cos\varphi\bigr]_0^\pi\cdot2\pi=\frac{R^3}3\cdot2\cdot2\pi=\frac43\pi R^3 .
\]
\end{example}

\begin{example}[Area of an ellipse]
The ellipse $x^2/a^2+y^2/b^2\le1$ is the image of the unit disc under $(x,y)=(au,bv)$, with $\abs{\det\Jac}=ab$. So $\Area=ab\cdot\pi\cdot1^2=\pi ab$.
\end{example}

\begin{example}[Moment of inertia of a disc]
For a disc of radius $R$ and constant density $\sigma$ about its axis: $I=\iint\sigma r^2\,\dd A=\sigma\int_0^{2\pi}\!\!\int_0^Rr^2\cdot r\,\dd r\,\dd\theta=\sigma\cdot2\pi\cdot\frac{R^4}4=\frac12(\sigma\pi R^2)R^2=\frac12mR^2$.
\end{example}

%====================================================================
\section{Improper integrals and the Gaussian integral}
%====================================================================

\begin{theorem}[Gaussian integral]\label{thm:gaussint}
$\displaystyle\int_{-\infty}^\infty e^{-x^2}\,\dd x=\sqrt\pi$.
\end{theorem}

\begin{proof}
Let $I_R=\int_{-R}^Re^{-x^2}\dd x$; we will show $I_R^2\to\pi$, which also proves that the improper integral $I=\lim_{R\to\infty}I_R$ exists (the numbers $I_R$ increase with $R$). We bound $I_R^2$ using discs and squares. Let $S_R=[-R,R]^2$, $D_R$ the disc of radius $R$. Since $D_R\subseteq S_R\subseteq D_{\sqrt2R}$ and the integrand $e^{-(x^2+y^2)}$ is positive,
\[
\iint_{D_R}e^{-(x^2+y^2)}\dd A\le\iint_{S_R}e^{-(x^2+y^2)}\dd A=I_R^2\le\iint_{D_{\sqrt2R}}e^{-(x^2+y^2)}\dd A .
\]
The middle equality is Fubini: $\iint_{S_R}e^{-x^2}e^{-y^2}=\bigl(\int_{-R}^Re^{-x^2}\bigr)^2$. In polar coordinates,
\[
\iint_{D_R}e^{-(x^2+y^2)}\dd A=\int_0^{2\pi}\!\!\int_0^Re^{-r^2}\,r\,\dd r\,\dd\theta=2\pi\cdot\tfrac12\bigl(1-e^{-R^2}\bigr)=\pi\bigl(1-e^{-R^2}\bigr).
\]
Both bounds tend to $\pi$ as $R\to\infty$, so $I^2=\lim I_R^2=\pi$. Since $I>0$, $I=\sqrt\pi$.
\end{proof}

\begin{intuition}
The function $e^{-x^2}$ has no elementary antiderivative, yet its total area is exactly $\sqrt\pi$. The trick: the \emph{square} of the integral is a \emph{two-dimensional} integral, and in polar coordinates the extra factor $r$ makes the integrand integrable by hand. This is a typical interview question.
\end{intuition}

%====================================================================
\section{Differentiation under the integral sign}
%====================================================================

\begin{theorem}[Leibniz rule]\label{thm:leibniz}
Let $f(x,t)$ and $\partial f/\partial t$ be continuous on $[a,b]\times[c,d]$, and put $F(t)=\int_a^bf(x,t)\,\dd x$. Then $F$ is $C^1$ on $[c,d]$ and
\[
F'(t)=\int_a^b\pd ft(x,t)\,\dd x .
\]
\end{theorem}

\begin{proof}
For $t$ and $t+h$ in $[c,d]$, the mean value theorem gives, for each $x$, a number $\theta=\theta(x,h)\in(0,1)$ with $f(x,t+h)-f(x,t)=h\,f_t(x,t+\theta h)$. Then
\[
\frac{F(t+h)-F(t)}h-\int_a^bf_t(x,t)\,\dd x=\int_a^b\bigl[f_t(x,t+\theta h)-f_t(x,t)\bigr]\dd x .
\]
Since $f_t$ is uniformly continuous on the compact rectangle, for any $\eps>0$ there is $\delta>0$ such that the bracket is smaller than $\eps$ in absolute value whenever $\abs h<\delta$. So the left side is at most $\eps(b-a)$ in absolute value, proving $F'(t)=\int f_t$. The function $t\mapsto\int f_t(x,t)\,\dd x$ is continuous by the same uniform continuity, so $F$ is $C^1$.
\end{proof}

\begin{corollary}[Variable limits]\label{cor:leibnizvar}
If moreover $\alpha,\beta:[c,d]\to[a,b]$ are $C^1$, then for $G(t)=\int_{\alpha(t)}^{\beta(t)}f(x,t)\,\dd x$,
\[
G'(t)=f(\beta(t),t)\,\beta'(t)-f(\alpha(t),t)\,\alpha'(t)+\int_{\alpha(t)}^{\beta(t)}\pd ft(x,t)\,\dd x .
\]
\end{corollary}

\begin{proof}
Let $H(u,v,t)=\int_u^vf(x,t)\,\dd x$. By the fundamental theorem of calculus $H_v=f(v,t)$ and $H_u=-f(u,t)$, and by Theorem~\ref{thm:leibniz} $H_t=\int_u^vf_t(x,t)\,\dd x$. All three partial derivatives are continuous, so $H$ is $C^1$ (Theorem~\ref{thm:c1diff}). Since $G(t)=H(\alpha(t),\beta(t),t)$, the chain rule gives $G'=H_u\alpha'+H_v\beta'+H_t$.
\end{proof}

\begin{example}[A classical trick]
Let $I(t)=\displaystyle\int_0^1\frac{x^t-1}{\ln x}\,\dd x$ for $t\ge0$. The integrand extends continuously to $x=0$ and $x=1$, and $\partial_t\frac{x^t-1}{\ln x}=x^t$ is continuous. By the Leibniz rule
\[
I'(t)=\int_0^1x^t\,\dd x=\frac1{t+1}.
\]
Since $I(0)=0$, we get $I(t)=\ln(1+t)$. A direct computation of this integral would be very hard.
\end{example}

\begin{pitfall}
The theorem needs hypotheses. If the limits are infinite, or $f_t$ is not continuous, exchanging $\frac{\dd}{\dd t}$ and $\int$ can be wrong. (In the infinite case one needs a dominating integrable function: this is the ``dominated convergence'' condition.)
\end{pitfall}

%====================================================================
\begin{summarybox}
\begin{itemize}
  \item $\iint_Rf\,\dd A$ is the limit of Riemann sums; continuous functions are integrable.
  \item Fubini: for continuous $f$, double and triple integrals equal iterated integrals in any order (proof by squeezing between lower and upper sums). Choose the order that makes the inner integral easy.
  \item Change of variables: $\dd x\,\dd y=\abs{\det\Jac\vec\varphi}\,\dd u\,\dd v$. Proof: small squares map to near-parallelograms; apply the inequality to $\vec\varphi$ and to $\vec\varphi^{-1}$.
  \item Polar $r\,\dd r\,\dd\theta$; cylindrical $r\,\dd r\,\dd\theta\,\dd z$; spherical $\rho^2\sin\varphi\,\dd\rho\,\dd\varphi\,\dd\theta$.
  \item $\int_{-\infty}^{\infty}e^{-x^2}\dd x=\sqrt\pi$ via polar coordinates in the plane.
  \item $\frac{\dd}{\dd t}\int_a^bf(x,t)\dd x=\int_a^b f_t\,\dd x$ when $f,f_t$ are continuous; with variable limits add the boundary terms.
\end{itemize}
\end{summarybox}

\begin{interview}
\begin{enumerate}
  \item \emph{Why is there an extra factor $r$ in polar coordinates?} The Jacobian determinant is $r$: a small coordinate rectangle $\dd r\times\dd\theta$ becomes a patch with sides $\dd r$ and $r\,\dd\theta$.
  \item \emph{Compute $\int_{-\infty}^\infty e^{-x^2}\dd x$.} Square it, use Fubini, switch to polar coordinates, get $\pi$.
  \item \emph{When can you swap the order of integration?} For continuous (or non-negative, or absolutely integrable) integrands. Show the counterexample with $(x^2-y^2)/(x^2+y^2)^2$.
  \item \emph{What is the volume element in spherical coordinates and why?} $\rho^2\sin\varphi\,\dd\rho\,\dd\varphi\,\dd\theta$: three perpendicular edges of lengths $\dd\rho$, $\rho\,\dd\varphi$, $\rho\sin\varphi\,\dd\theta$.
  \item \emph{What does $\det\Jac$ mean when it is negative?} The map reverses orientation; the integral uses $\abs{\det\Jac}$.
  \item \emph{Differentiate $\int_0^{t}f(x,t)\dd x$ with respect to $t$.} $f(t,t)+\int_0^tf_t\,\dd x$.
\end{enumerate}
\end{interview}

\subsection*{Exercises}
\begin{exercise}
Compute $\displaystyle\int_0^2\!\!\int_{y/2}^1\cos(x^2)\,\dd x\,\dd y$ by swapping the order. (The region is $0\le x\le1$, $0\le y\le2x$. Answer: $\sin1$.)
\end{exercise}
\begin{exercise}
Find the volume of the region inside both $x^2+y^2+z^2\le4$ and $x^2+y^2\le1$, using cylindrical coordinates. (Answer: $\tfrac{4\pi}3(8-3\sqrt3)$.)
\end{exercise}
\begin{exercise}
Use $u=x+y$, $v=x-y$ to compute $\iint_D(x+y)^2\dd A$ over the square $\abs x+\abs y\le1$. (Hint: $\abs x+\abs y=\max(\abs u,\abs v)$, so the region becomes $\abs u\le1,\abs v\le1$, and $\dd x\,\dd y=\tfrac12\dd u\,\dd v$. Answer: $\tfrac23$.)
\end{exercise}
\begin{exercise}
Show that $\int_0^\infty e^{-x^2}\dd x=\frac{\sqrt\pi}2$ and deduce $\Gamma(\tfrac12)=\sqrt\pi$, where $\Gamma(s)=\int_0^\infty x^{s-1}e^{-x}\dd x$ (substitute $x=u^2$).
\end{exercise}
\begin{exercise}
Compute the volume of the $n$-dimensional unit ball for $n=2,3$ by Fubini, and find a recursion $V_n=V_{n-1}\int_{-1}^1(1-t^2)^{(n-1)/2}\dd t$.
\end{exercise}
\begin{exercise}
Using Corollary~\ref{cor:leibnizvar}, compute $\dfrac{\dd}{\dd t}\displaystyle\int_t^{t^2}\frac{\sin(tx)}x\,\dd x$. (Answer: $\bigl(3\sin t^3-2\sin t^2\bigr)/t$.)
\end{exercise}
